\clearpage\begingroup\SourceLeaf{394b}\ArithmeticHead{63}{R}{DIE KOMPLEXEN GRÖSSEN IN IHRER BEZIEHUNG ZUR GEOMETRIE.}
\renewcommand{\thefootnote}{\arabic{footnote})}\setcounter{footnote}{0}
\begin{center}
\textbf{17.}\par\smallskip
\textls[100]{Biquadratische Reste.}
\end{center}
\noindent\hspace*{\parindent}Als \textsc{Gauss} im Jahre 1805 von den quadratischen Resten zu den kubischen
und biquadratischen fortschritt, fand er sogleich durch Induktion eine Reihe
einfacher Lehrsätze, die mit den für die quadratischen Reste gefundenen Ergebnissen
überraschende Ähnlichkeit hatten, jedoch ist es ihm erst nach vielen,
durch eine Reihe von Jahren fortgesetzten Versuchen gelungen, befriedigende
Beweise dafür aufzufinden. Zu diesem Zwecke musste er neue Wege einschlagen,
nämlich »das Feld der höhern Arithmetik, welches man sonst nur
auf die reellen ganzen Zahlen ausdehnte, auch über die imaginären erstrecken
und diesen das völlig gleiche Bürgerrecht mit jenen einräumen« (W. II, S. 171).
Wie es scheint ist diese »erlösende Eingebung« in das Jahr 1807 zu setzen
(\textsc{Bachmann}, W. X\textsubscript{2}, Abh. 1, S. 55). Vollständig durchgedrungen ist \textsc{Gauss} freilich
erst 1813 (T. Nr. 144) und veröffentlicht hat er seine Untersuchungen erst 1831
in der Abhandlung über die biquadratischen Reste (W. II, S. 93), die er durch
die wiederholt erwähnte Selbstanzeige noch ergänzte (W. II, S. 169). »Wie
einfach jetzt auch eine solche Einführung der komplexen Zahlen als Moduln
erscheinen mag«, hat \textsc{Jacobi}\footnote{C. G. J. \textsc{Jacobi}, \textit{Über die complexen Primzahlen}, \textsc{Crelles} Journal, Bd. 19, 1839, S. 314, Werke VI, S. 275.} geurteilt, »so gehört sie nichtsdestoweniger zu
den tiefsten Gedanken der Wissenschaft; ja ich glaube nicht, dass zu einem
so verborgenen Gedanken die Arithmetik allein geführt hat, sondern dass er
aus dem Studium der elliptischen Transzendenten geschöpft ist, und zwar aus
der besonderen Gattung derselben, welche die Rektifikation von Bogen der
Lemniscata gibt. In der Theorie der Vervielfachung und Teilung von Bogen
der Lemniscata spielen nämlich die komplexen Zahlen von der Form \(a+bi\)
genau die Rolle gewöhnlicher Zahlen. . . . So wie man einen Kreisbogen,
wenn man ihn in 15 Teile teilen soll, in 3 und in 5 Teile teilt und aus beiden
\clearpage\SourceLeaf{395a}\ArithmeticHead{64}{L}{STÄCKEL, GAUSS ALS GEOMETER.}\setcounter{footnote}{0}
\noindent Teilungen die gesuchte findet, so hat man einen Bogen der Lemniscata, um
ihn in 17 Teile zu teilen, in \(1+4i\) und \(1-4i\) Teile zu teilen, und setzt die
Teilung in 17 Teile aus beiden zusammen«.
\par
Ebenso wichtig wie diese Erweiterung des Zahlengebietes, mit der die
Lehre von den algebraischen Zahlen ins Leben gerufen wurde, ist für die
Fortschritte der höheren Arithmetik die Darstellung der ganzen komplexen
Zahlen vermöge der Gitterpunkte der Ebene geworden. Hieran schliesst sich
bei der Untersuchung der ternären quadratischen Form die Heranziehung der
Gitterpunkte im Raume (W. II, S. 188). Es ist sogar wahrscheinlich, dass
\textsc{Gauss} bereits Zahlengitter im Raume von \(n\) Dimensionen betrachtet hat, denn
die Andeutung nach dieser Richtung, die \textsc{Eisenstein} 1844 gemacht hat, geht
wohl auf seinen Aufenthalt in Göttingen während des Sommers dieses Jahres
zurück\footnote{\textsc{Eisenstein}, \textit{Geometrischer Beweis des Fundamentaltheorems für die quadratischen Reste}, \textsc{Crelles} Journal, Bd. 28, 1844, S. 248.}. So muss \textsc{Gauss} auch als der Begründer der Geometrie der Zahlen
gelten.
\par
Die Bedeutung der beiden Veröffentlichungen vom Jahre 1831 geht jedoch
über die \textls[100]{Zahlentheorie} hinaus. Wenn \textsc{Gauss} in der neuen Darstellung des
Beweises für den Fundamentalsatz der Algebra, den er 1849 gab, sagt: »gegenwärtig,
wo der Begriff der \textls[100]{komplexen Grössen} jedermann geläufig ist« (W.
III, S. 74), so hat die Analysis ihm diesen Fortschritt zu verdanken. Gewiss
hatten schon \textsc{Wessel} (1799), \textsc{Argand} (1806) und andere nach ihnen die selbständige
Berechtigung und die geometrische Darstellung der komplexen Grössen
erkannt und wichtige Anwendungen davon zu machen gewusst, allein die
Kenntnis und Würdigung ihrer Untersuchungen ist auf enge Kreise beschränkt
geblieben. Es bedurfte eines \textsc{Gauss}, um die Hemmungen zu beseitigen und
die neuen Anschauungen zum Siege zu führen.
\clearpage\endgroup

